\documentclass{psutex}
\usepackage[utf8]{inputenc}
\usepackage[T1]{fontenc}
\usepackage{geometry}
\usepackage{graphicx}
\usepackage{float}
\usepackage{hyperref}
\usepackage[english,indonesian]{babel}

\title{Optimasi Pemrosesan Template LaTeX Menggunakan Pendekatan Berbasis AI}
\author{Academic Support Agent}
\date{\today}

\begin{document}
\maketitle

\begin{abstract}

\end{abstract}

\chapter{PENDAHULUAN}


\chapter{TINJAUAN PUSTAKA}


% Original content (commented out):
% \documentclass[double,12pt]{psutex}
\usepackage{graphicx}
\usepackage{rotating} %Package added to allow the rotation of figures and chart on a page, {sidewaysfigure} command
\usepackage{tablefootnote} %Packaged added to allow footnotes in the tabular environment, use \tablefootnote command

\usepackage{url}

\hyphenation{op-tical net-works semi-conduc-tor}

\usepackage{graphicx}
\usepackage{amsmath}
\usepackage{amssymb}
\usepackage{gensymb} % for degree symbol
\usepackage[dvipsnames]{xcolor} % colorbox for inline comments
\usepackage{cite}
\usepackage{booktabs}
\usepackage{tabularx}
\usepackage[flushleft]{threeparttable}
\usepackage{dcolumn}
\usepackage{float}
\usepackage{longtable}

\DeclareMathOperator*{\argmin}{argmin}
\DeclareMathOperator*{\argmax}{argmax}
\newcommand{\figscale}{1.0}

\newcommand{\shortnote}[1]{\footnote{\color{red}#1}}

\renewcommand{\arraystretch}{1.0}


% -------------------------------------------------
% Define variables needed to build the Title Page
% -------------------------------------------------
\title{A One or Two Line Title of the Thesis \\ Second Line as Needed} % line break needs to be inserted manually for it to be double-spaced
\author{Some Student}
\graddegree{Master of Science}
\doctype{Thesis}
\university{Portland State University}
\department{Electrical and Computer Engineering}
\major{Electrical and Computer Engineering}
\submitdate{November, 2022}
\commencementyear{2022}
\committeemembers{Thesis Advisor, Chair\\
Other CommitteeMember\\
Another CommitteeMember}

\usepackage{etoolbox}           % For the flag determining if front matter goes into the TOC
\newtoggle{fulltoc}
\toggletrue{fulltoc}  % Change to \togglefalse{fulltoc} to remove front matter



\begin{document}

\maketitle

%-------------------------Pretext Pages----------------------------
% Abstract
\include{chapters/02-abstract}

% Dedication
\include{chapters/04-dedication}

% Acknowledgements
\include{chapters/06-ack}

%-------------------------Listings-----------------------------
% Table of contents, list of figures, etc
% Manually uncomment the listings for types that are present in the document
\tableofcontents
\listoftables
\listoffigures
% \listofappendices
% \listofappendixfigures
% \listofappendixtables

%-------------------------Main Chapters-----------------------------
\mainmatter

\include{chapters/10-intro}
\include{chapters/40-chap2.tex}
\include{chapters/60-results}
\include{chapters/80-conclusion}
\include{chapters/90-bib}
\end{document}

\end{document}